\uuid{msxR}
\titre{Quantile d'une loi normale}
\auteur{Maxime Nguyen}
\organisation{AMSCC}
\datecreate{2026-11-27}
\niveau{L2}
\module{Probabilité}
\chapitre{Probabilité continue}
\sousChapitre{Loi normale}
\difficulte{1}
\theme{loi normale, quantile, standardisation}

\contenu{
\begin{question}
\qtype{onlyone}
On note $\mathcal{N}(\mu,\sigma)$ la loi normale d'espérance $\mu$ et d'écart-type $\sigma$, et $\Phi$ la fonction de répartition de $\mathcal{N}(0,1)$. On donne $\Phi(1{,}96) \approx 0{,}975$. Soit $X \sim \mathcal{N}(20, 4)$. Le réel $a$ tel que $\mathbb{P}(X \leq a) = 0{,}975$ vaut environ :
\begin{answers}
\bad{$1{,}96$}
\bad{$12{,}16$}
\good{$27{,}84$}
\bad{$51{,}36$}
\end{answers}
\begin{explanations}
$\mathbb{P}(X \leq a) = \Phi\left(\frac{a - 20}{4}\right) = 0{,}975$
donne $\frac{a - 20}{4} = 1{,}96$, soit $a = 20 + 4 \times 1{,}96 = 27{,}84$.

A est le quantile de $\mathcal{N}(0,1)$ : on a oublié de revenir à $X$.

B : $20 - 4 \times 1{,}96$, erreur de signe ; on sait que $a > 20$ puisque la probabilité dépasse $\frac12$.

D : $20 + 16 \times 1{,}96$, la variance a pris la place de l'écart-type.

Contre-sens visés : s'arrêter au quantile de la loi centrée réduite ;
multiplier par la variance. Contrôle rapide : une probabilité supérieure à
$\frac12$ donne un quantile supérieur à $\mu$.
\end{explanations}
\end{question}
}
